% -----------------------------------------------------------------------------------
\section{Decision problems with prediction}\label{sec:model}

This section fixes the class of problems on which decision theories are compared. The aim is a
model rich enough to contain the problems that drive the debate (Newcomb's problem, the twin
prisoner's dilemma, Death in Damascus, the Bomb, counterfactual mugging) and precise enough that
every criterion to be considered has a determinate verdict on every problem. Each place where an
informal notion---a counterfactual, \enquote{the same computation}, \enquote{holding the state
fixed}---has to be settled by stipulation is marked as such; the stipulations are collected in
Appendix~\ref{app:stipulations}.

\subsection{Problems, policies, and structural plans}

\begin{definition}[Decision problem]\label{def:problem}
A \emph{decision problem} $D$ is a finite extensive-form tree with three kinds of non-terminal
nodes and real payoffs at its leaves:
\begin{enumerate}[label=(\roman*)]
  \item \emph{nature nodes}, with fixed probability distributions over their successors;
  \item \emph{prediction nodes}, nature nodes whose distribution may depend on the agent's policy,
        but only through the actions that the policy prescribes at specified information sets;
  \item \emph{agent nodes}, partitioned into a finite set $H_D$ of information sets, each $h\in H_D$
        carrying a finite nonempty set $A_h$ of actions.
\end{enumerate}
It is required that no information set occurs twice on any path of the tree.
\end{definition}

Perfect recall implies the last requirement, which is all that will be used. A \emph{deterministic
policy} for $D$ is a choice of action at every information set, and the policy space is
\[
  \Pol_D=\prod_{h\in H_D}A_h,
\]
a finite nonempty set. Mixed policies are deferred to Sections~\ref{sec:robust}
and~\ref{sec:mixed}.

The policy-dependent kernels of the prediction nodes determine the distribution of actual play,
but they do not by themselves settle the causal counterfactuals that several criteria require:
what would have happened had the agent acted otherwise, \emph{with the environment's responses
held fixed}. To make these counterfactuals explicit, the tree is supplemented with a structural
description of the environment.

\begin{definition}[Structural representation]\label{def:structural}
A \emph{structural representation} of $D$ consists of
\begin{enumerate}[label=(\roman*)]
  \item a finite nonempty set $S_D$ of \emph{environment plans}, each of which specifies the response
        of every nature and prediction node, including nodes on branches not taken and responses that
        depend on earlier physical actions;
  \item for every policy $\pi\in\Pol_D$, a probability distribution $\mu_\pi\in\Delta(S_D)$;
  \item the payoff $U_D(s,\pi)$ obtained by traversing the tree with plan $s$ and policy $\pi$.
\end{enumerate}
The \emph{ex-ante value} of $\pi$ is
\[
  V_D(\pi)=\sum_{s\in S_D}\mu_\pi(s)\,U_D(s,\pi).
\]
\end{definition}

The representation separates two kinds of dependence. \emph{Ordinary causal dependence} of the
environment on earlier actions is carried by the responses written into each plan $s$.
\emph{Predictive dependence} on the policy is carried by the map $\pi\mapsto\mu_\pi$. A tree with
policy-dependent kernels can always be given a structural representation by supplying potential
outcomes for its nature and prediction nodes; how the outcomes on alternative branches are coupled
is an additional modelling choice, which matters for the causal criteria and is fixed by the
representation. Conversely, every structural representation is implemented by a tree whose root
is a prediction node selecting $s$ according to $\mu_\pi$, followed by deterministic traversal.

\subsection{Well-definedness and fairness}

Two conditions on problems are distinguished.

\begin{definition}[Well-definedness]\label{def:welldefined}
$D$ is \emph{well defined} if $\mu_\pi$ is a probability distribution on $S_D$ for \emph{every}
$\pi\in\Pol_D$.
\end{definition}

\begin{definition}[Fairness]\label{def:fair}
$D$ is \emph{fair} if $S_D$, $U_D$, and the entire map $\pi\mapsto\mu_\pi$ are fixed features of $D$:
the evaluation of an agent may depend on the actions its policy prescribes, but not on the
identity, name, or source code of the decision theory that produced the policy. Fixed algorithms
may occur as constants of the environment (\enquote{the predictor simulates the algorithm $g_0$}),
provided $g_0$ does not vary with the agent being evaluated.
\end{definition}

Fairness excludes the trivial diagonal problem that pays an agent exactly when its action differs
from that of a given theory $T$; such a problem defeats every $T$ and shows nothing about decision
theory. Well-definedness is a separate requirement. A verbally described \enquote{perfect
predictor} may satisfy an action-based fairness condition while its promised behaviour is
logically inconsistent for some policy; Section~\ref{sec:bomb} gives an example. The presence of a
fixed algorithm in the environment does not by itself make a self-referential specification well
defined (Remark~\ref{rem:skalse}). Fixed algorithms are harmless here because a problem is a finite
object whose data include their outputs; when problems are given as source code and the agent must
compute those outputs, they are not (Section~\ref{sec:computation}).

Write $\cls^*$ for the class of well-defined, fair finite problems with a fixed structural
representation, closed under the commitment construction of Section~\ref{sec:commit}.

\begin{definition}[Decision theory]\label{def:theory}
A \emph{decision theory} on a class $\cls$ of problems assigns to every $D\in\cls$ a policy
$T_D\in\Pol_D$. No uniformity, computability, or generality is required: every lookup table is a
decision theory.
\end{definition}

Totality is part of the definition. A proposed criterion that recommends nothing on some problem
therefore fails to define a decision theory on any class containing that problem; it does not
define a theory that is silent there. The permissive definition is the weakest sense of
\enquote{decision theory}, and it is the first of three. Part~\ref{part:across} requires a theory to be a
single rule on a common reduced form of problems, and Part~\ref{part:programs} requires it to be a
computable procedure that receives a description of the problem and must itself work out what the
structural representation supplies here.

\subsection{Reach and local causal deviation}

For a policy $\pi$ and an information set $h$, let
\[
  R_h(\pi)=\{s\in S_D: h\text{ is reached in the traversal with }(s,\pi)\},
  \qquad
  r_\pi(h)=\mu_\pi\bigl(R_h(\pi)\bigr).
\]
The number $r_\pi(h)$ is the actual probability that $h$ is reached under $\pi$; $h$ is
\emph{positively reached} under $\pi$ if $r_\pi(h)>0$. For $a\in A_h$, let $\dev{\pi}{h}{a}$ denote the
policy that agrees with $\pi$ everywhere except at $h$, where it prescribes $a$.

\begin{lemma}\label{lem:local}
For every $\pi$, $h$, and $a\in A_h$:
\begin{enumerate}[label=(\roman*)]
  \item $R_h(\dev{\pi}{h}{a})=R_h(\pi)$;
  \item if $s\notin R_h(\pi)$, then $U_D(s,\dev{\pi}{h}{a})=U_D(s,\pi)$.
\end{enumerate}
\end{lemma}

\begin{proof}
With the plan $s$ fixed, the path through the tree is determined by the prescriptions at the
information sets it visits. Up to its first visit to $h$, the path under $(s,\dev{\pi}{h}{a})$ uses
only prescriptions at information sets other than $h$, which are unchanged; so it coincides with the
path under $(s,\pi)$, and $h$ is reached under one exactly when it is reached under the other. If
$h$ is not reached, the two paths coincide entirely. Since no information set recurs on a path, a
visit to $h$ is the first and only one.
\end{proof}

The \emph{causal value} of $a$ at a positively reached $h$, with beliefs induced by $\pi$, is
\[
  Q_\pi(h,a)=\frac{1}{r_\pi(h)}\sum_{s\in R_h(\pi)}\mu_\pi(s)\,U_D\bigl(s,\dev{\pi}{h}{a}\bigr).
\]
This fixes the causal counterfactual completely: the distribution over environmental responses is
the one the actual policy induces, the action at $h$ is changed, and the policy's other
prescriptions are kept. The contrast with ex-ante evaluation lies in a single subscript:
$Q_\pi(h,a)$ averages with $\mu_\pi$, whereas $V_D(\dev{\pi}{h}{a})$ averages with
$\mu_{\dev{\pi}{h}{a}}$. Every within-problem conflict studied below is a consequence of that
difference.
